\subsection{不可约表示}

引理 1. --- G 的每个不可约酉表示都是平方可积的。

事实上，不可约酉表示的矩阵系数连续且有界，所以在 G 上平方可积，因为 G 是紧的。

命题 2. --- 设 π 是 G 在希尔伯特空间 E 上的不可约酉表示。E 的维数有限，且等于 π 的形式次数。

表示 \(\pi\) 是平方可积的（引理 1），所以形式次数 \(c_{\mu}(\pi)\) 相对于 \(\mu\) 有定义（V, p. 423, 定义 4）；它是严格为正的实数。设 \((e_i)_{i \in I}\) 是 E 中的有限标准正交族。设 \(x\) 是 E 中范数为 1 的元（这样的元存在，因为 E 非零）。对 \(i \in I\)，有

\[
c _ {\mu} (\pi) \int_ {\mathrm{G}} | \langle e _ {i} | \pi (g) x \rangle | ^ {2} d \mu (g) = 1
\]

（V, p. 424, 命题 8）。对 \(i \in I\) 将此公式求和，得到

\[
\begin{array}{r l} & {\mathrm{Card(I)} = c _ {\mu} (\pi) \int_ {\mathrm{G}} \sum_ {i \in \mathrm{I}} | \langle e _ {i} | \pi (g) x \rangle | ^ {2} d \mu (g)} \\ & {\qquad \leqslant c _ {\mu} (\pi) \int_ {\mathrm{G}} \| \pi (g) x \| ^ {2} d \mu (g) = c _ {\mu} (\pi)} \end{array}
\]

由贝塞尔不等式（EVT, V, p. 21, 命题 4）。因此 I 的基数上界为 \(c_{\mu}(\pi)\)。这意味着 E 的维数有限。

于是可将前述结果应用于 E 的标准正交基 \((e_{i})_{i\in\mathrm{I}}\)。由 EVT, V, p. 22, 命题 5，得到

\[
\begin{array}{r} \dim (\mathrm{E}) = c _ {\mu} (\pi) \int_ {\mathrm{G}} \sum_ {i \in \mathrm{I}} | \langle e _ {i} | \pi (g) x \rangle | ^ {2} d \mu (g) \\ = c _ {\mu} (\pi) \int_ {\mathrm{G}} \| \pi (g) x \| ^ {2} d \mu (g) = c _ {\mu} (\pi), \end{array}
\]

证明完毕。

特别地，因此可以谈论特征标 \(\chi_{\pi}\)，其中 \(\pi\) 是 G 的一个不可约酉表示。它是 G 上的连续函数。

推论 1. --- 设 \(\pi\) 是 G 在希尔伯特空间 E 上的不可约酉表示。\(\pi\) 的特征标是对合代数 \(\mathrm{L}^{1}(\mathrm{G})\) 的埃尔米特元。

设 \(x \in G\)。由于 \(G\) 是单模的，有 \(\chi_{\pi}^{*}(x) = \overline{\operatorname{Tr}(\pi(x^{-1}))}\)，又有 \(\chi_{\pi}^{*}(x) = \operatorname{Tr}(\pi(x))\)，这是由于 \(\pi\) 是酉表示。

推论 2. --- a) 设 \(\pi\) 是 G 在希尔伯特空间 E 上的不可约酉表示。有

\[
\int_ {\mathrm{G}} \overline {{\langle x \mid \pi (g) x ^ {\prime} \rangle}} \langle y \mid \pi (g) y ^ {\prime} \rangle d \mu (g) = \frac {1}{\dim (\mathrm{E})} \overline {{\langle x \mid y \rangle}} \langle x ^ {\prime} \mid y ^ {\prime} \rangle
\]

对所有 \((x,y,x',y')\in \mathrm{E}^4\) 都成立。

\begin{enumerate}
\def\labelenumi{\alph{enumi})}
\setcounter{enumi}{1}
\tightlist
\item
  设 \(\pi_{1}\)（分别为 \(\pi_{2}\)）是 G 在希尔伯特空间 \(E_{1}\)（分别为 \(E_{2}\)）上的不可约酉表示。若 \(\pi_{1}\) 和 \(\pi_{2}\) 不同构，则有
\end{enumerate}

\[
\int_ {G} \overline {{\langle x \mid \pi_ {1} (g) x ^ {\prime} \rangle}} \langle y \mid \pi_ {2} (g) y ^ {\prime} \rangle d \mu (g) = 0
\]

对所有 \((x,x^{\prime},y,y^{\prime})\in \mathrm{E}_{1}^{2}\times \mathrm{E}_{2}^{2}\) 都成立

这立即由 V, p. 424, 命题 8 和 V, p. 424, 命题 9 得出，同时考虑引理 1 及公式 \(c_{\mu}(\pi) = \dim(\mathrm{E})\)（命题 2）。

推论 3. --- 设 \(\pi_1\) 和 \(\pi_2\) 是 G 的不可约表示。在希尔伯特空间 \(\mathrm{L}^{2}(\mathrm{G})\) 中，有 \(\langle \chi_{\pi_1} | \chi_{\pi_2} \rangle = 1\)（当 \(\pi_1\) 和 \(\pi_2\) 同构时），否则有 \(\langle \chi_{\pi_1} | \chi_{\pi_2} \rangle = 0\)。换言之，类 \(\pi \in \widehat{\mathbf{G}}\) 的特征标所成的族是 \(\mathrm{L}^{2}(\mathrm{G})\) 中的标准正交族。

这由 V, p. 457, 命题 1 及推论 2 得出；注意，对于 G 的有限维酉表示的空间的每个标准正交基 \((e_{i})_{i\in\mathrm{I}}\)，以及表示 \(\pi\) 和每个 \(g\in G\)，都有公式

\[
\chi_ {\pi} (g) = \sum_ {i \in I} \langle e _ {i} | \pi (g) e _ {i} \rangle .
\]

推论 4. --- a) 设 \(\varrho\) 是 G 的连续有限维表示。有

\[
\chi_ {\varrho} = \sum_ {\pi \in \widehat {G}} \dim (\operatorname{Hom} _ {G} (\pi , \varrho)) \chi_ {\pi} = \sum_ {\pi \in \widehat {G}} \dim (\operatorname{Hom} _ {G} (\varrho , \pi)) \chi_ {\pi}.
\]

\begin{enumerate}
\def\labelenumi{\alph{enumi})}
\setcounter{enumi}{1}
\tightlist
\item
  设 \(\varrho_{1}\) 和 \(\varrho_{2}\) 是 G 的连续有限维表示。有
\end{enumerate}

\[
\left\langle \chi_ {\varrho_ {1}} \mid \chi_ {\varrho_ {2}} \right\rangle = \dim (\operatorname{Hom} _ {\mathrm{G}} (\varrho_ {1}, \varrho_ {2})) = \dim (\operatorname{Hom} _ {\mathrm{G}} (\varrho_ {2}, \varrho_ {1})).
\]

\begin{enumerate}
\def\labelenumi{\alph{enumi})}
\setcounter{enumi}{2}
\tightlist
\item
  G 的连续有限维表示 \(\varrho\) 不可约，当且仅当 \(\| \chi_{\varrho}\|^{2} = 1\)。
\end{enumerate}

由于 \(\varrho\) 是半单的，它同构于其 \(\pi\)-同型分量 \(\mathrm{M}_{\pi}(\varrho)\) 的直和，其中 \(\pi \in \widehat{\mathrm{G}}\)。以 \(m_{\pi}(\varrho)\) 表示 \(\pi\) 在 \(\varrho\) 中的重数（V, p. 398, 定义 10）；于是有

\[
\chi_ {\varrho} = \sum_ {\pi \in \widehat {G}} m _ {\pi} (\varrho) \chi_ {\pi}.
\]

因此，断言 a) 由以下事实得出：

\[
m _ {\pi} (\varrho) = \dim \operatorname{Hom} _ {\mathrm{G}} (\pi , \varrho) = \dim \operatorname{Hom} _ {\mathrm{G}} (\varrho , \pi)
\]

（V, p. 377, 公式 (1) 及 V, p. 387, 推论 2）。

由特征标的双线性和标准正交性，断言 \(a\) ) 蕴含

\[
\langle \chi_ {\varrho_ {1}} \mid \chi_ {\varrho_ {2}} \rangle = \sum_ {\pi \in \widehat {\mathrm{G}}} \dim (\operatorname{Hom} _ {\mathrm{G}} (\pi , \varrho_ {1})) \dim (\operatorname{Hom} _ {\mathrm{G}} (\pi , \varrho_ {2})),
\]

另一方面，有典范同构

\[
\begin{array}{c} \operatorname{Hom} _ {\mathrm{G}} (\varrho _ {1}, \varrho _ {2}) \to \bigoplus_ {(\pi _ {1}, \pi _ {2}) \in \widehat {\mathrm{G}} \times \widehat {\mathrm{G}}} \operatorname{Hom} _ {\mathrm{G}} (\mathrm{M} _ {\pi _ {1}} (\varrho _ {1}), \mathrm{M} _ {\pi _ {2}} (\varrho _ {2})) \\ \qquad \qquad \qquad \qquad \qquad \qquad \qquad \qquad \qquad \qquad \qquad \qquad \qquad \qquad \qquad \qquad \qquad \qquad \qquad \qquad \qquad \qquad \qquad \qquad \qquad \qquad \qquad \qquad \qquad \qquad \qquad \qquad \qquad \qquad \qquad \to \bigoplus_ {\pi \in \widehat {\mathrm{G}}} \operatorname{Hom} _ {\mathrm{G}} (\mathrm{M} _ {\pi} (\varrho _ {1}), \mathrm{M} _ {\pi} (\varrho _ {2})) \end{array}
\]

（V, p. 377, 公式 (1)），从而可得

\[
\dim (\operatorname{Hom} _ {\mathrm{G}} (\varrho_ {1}, \varrho_ {2})) = \sum_ {\pi \in \widehat {\mathrm{G}}} m _ {\pi} (\varrho_ {1}) m _ {\pi} (\varrho_ {2}),
\]

从而得到断言 \(b\)。断言 \(c\) 也由 \(a\) 和舒尔引理（V, p. 386, 命题 6）得出。

推论 5. --- 设 \(\pi_1\) 和 \(\pi_2\) 是 G 的不可约酉表示。则 \(\pi_1(\overline{\chi}_{\pi_2}) = 0\)；若 \(\pi_1\) 不同构于 \(\pi_2\)，而 \(\pi_1(\overline{\chi}_{\pi_1})\) 是乘以 \(1/\dim(\pi_1)\) 的映射。

\(\pi_2\) 的特征标是 G 上的连续中心函数，因此线性映射 \(u = \pi_1(\overline{\chi}_{\pi_2})\) 有定义且属于空间 \(\mathrm{Hom}_{\mathrm{G}}(\pi_1,\pi_1)\)。由舒尔引理（V, p. 386, 命题 6），它是一个倍乘映射，其迹为

\[
\begin{array}{c} \operatorname{Tr} (u) = \operatorname{Tr} \Bigl (\int_ {\mathrm{G}} \overline {{\chi_ {\pi_ {2}} (g)}}   \pi_ {1} (g) d \mu (g) \Bigr) \\ = \int_ {\mathrm{G}} \overline {{\chi_ {\pi_ {2}} (g)}} \chi_ {\pi_ {1}} (g) d \mu (g). \end{array}
\]

由正交关系，\(u\) 的迹为零；若 \(\pi_1\) 不同构于 \(\pi_2\)，否则等于 1。断言得证。

注。--- 当 G 有限时，舒尔正交关系（分别为特征标正交公式）与 A, VIII, p. 399 中的公式（分别与 A, VIII, p. 400, 命题 4 中的公式）一致。另一方面，有限群特征标的“第二正交关系”（A, VIII, p. 402, 公式 (32)）在 G 紧时没有完全对应的形式。

当 G 有限时，与 e 的共轭类对应的第二正交公式的特殊情形是公式

\[
\frac {1}{\mathrm{CardG}} \sum_ {\pi \in \widehat {\mathrm{G}}} \dim (\pi) \chi _ {\pi} (g) = \left\{ \begin{array}{l l} 1 & \text {if} g = e \\ 0 & \text {otherwise}, \end{array} \right.
\]

这也可解释为计算 G 的正则表示 \(\gamma_{G}\) 的特征标，并且等价于对从 G 到 C 的每个函数 f 都有公式 \(\mathrm{Tr}(\boldsymbol{\gamma}_{\mathrm{G}}(f)) = f(e)\)。

再设 G 是任意紧群。对每个函数 \(f \in \mathcal{C}(\mathrm{G})\)，\(\boldsymbol{\gamma}_{\mathrm{G}}(f)\) 是 \(\mathrm{L}^{2}(\mathrm{G})\) 的自同态，与自同态 \(\varphi \mapsto f * \varphi\) 重合，且具有有限迹（V, p. 407, 引理 4 及 III, p. 33, 推论 2）。由同上及 IV, p. 177, 定理 2，还有

\[
\operatorname{Tr} \left(\boldsymbol {\gamma} _ {\mathrm{G}} (f)\right) = \int_ {\mathrm{G}} f \left(x ^ {- 1} x\right) d \mu (x) = f (e).
\]

命题 3. --- 设 \(G = G_{1} \times G_{2}\)，其中 \(G_{1}\) 和 \(G_{2}\) 是紧群。从 \(\widehat{G}_{1} \times \widehat{G}_{2}\) 到 \(\widehat{G}\) 的映射 b 将 \((\pi_{1}, \pi_{2})\) 送到 \(\pi_{1} \boxtimes \pi_{2}\) 的类，且 b 是双射。

由 V, p. 389, 推论 6，映射 b 定义良好。

设 \(\pi_1\) 和 \(\pi_2\) 是 \(\widehat{\mathbf{G}}\) 的元素。设 \(\pi \in \widehat{\mathbf{G}}\)。\(\pi\)-同型分量是限制 \(\varpi\) 的 \(\pi_1 \boxtimes \pi_2\) 在 \(\mathrm{G}_1 \times \{e\}\) 上，其等于 \(\varpi\) 当 \(\pi_1 = \pi\) 时，否则为零（V, p. 384, 引理 8）。因此，酉表示 \(\pi_1 \boxtimes \pi_2\) 在同构意义下确定 \(\pi_1\)；同样，它确定 \(\pi_2\)，这证明 \(b\) 是单射。

证明 \(b\) 是满射。设 \(\pi\) 是 \(G_1 \times G_2\) 在希尔伯特空间 E 上的不可约酉表示。它是有限维的（命题 2）。设 \(\pi_1\) 是 \(G_1\) 在希尔伯特空间 \(E_1\) 上的不可约酉表示，使得 \(\pi\) 到 \(G_1 \times \{e\}\) 的限制包含一个与 \(\pi_1\) 同构的子表示（V, p. 379, 引理 3）。令 \(F = \text{Hom}_{G_1}(\pi_1, \pi)\)。这是非零有限维向量空间。对 \(h \in G_2\) 和 \(u \in F\)，令 \(\varrho(h)(u)\) 为从 \(E_1\) 到 E 的线性映射，由 \(x \mapsto \pi(e, h)(u(x))\) 定义。由于 \(G_1 \times \{e\}\) 与 \(\{e\} \times G_2\) 在 G 中交换，映射 \(\varrho(h)(u)\) 属于空间 F。映射 \(\varrho\) 是 \(G_2\) 在空间 F 上的线性表示；它是连续的。由于 F 非零，存在不可约子表示 \(\pi_2\) 的 \(\varrho\)（V, p. 379, 引理 3）。令 \(E_2\) 为 \(\pi_2\) 的空间。于是，线性映射 \(v: E_1 \otimes E_2 \to E\) 满足 \(x \otimes u \mapsto u(x)\) 对 \(x \in E_1\) 和 \(u \in E_2\) 成立，是从 \(\pi_1 \boxtimes \pi_2\) 到 \(\pi\) 的非零 G-态射；由于表示 \(\pi\) 和 \(\pi_1 \boxtimes \pi_2\) 都不可约且有限维，态射 \(v\) 是同构（V, p. 378, 引理 2）。

本命题可与 A, VIII, p. 208 的定理 1 比较。

